\chapter*{Preface}
\fontsize{12}{14.5}\selectfont

In this thesis I explore the notion of ideally exact category, introduced by George Janelidze in 2024. The main goal is to provide an overview of the ideally exact categories with the necessary background needed to understand it, and present some of the most important results in this field.

A basic knowledge of category theory — such as limits, colimits, adjunctions, and monads — is assumed throughout this thesis. Some familiarity with Malt'sev, exact, regular, protomodular and semiabelian categories would provide a natural foundation; however the main constructions and definitions of these categories are reviewed in the first chapter. The reader who find it helpful, can also refer to the appendix, where some results about these categories — necessary for proving several results in the thesis — are presented.

Right after the introduction and preliminaries, I will present in the second chapter the pullback functor and its left adjoint, which are a key tool in the theory of ideally exact categories. In this chapter I also define the connected limits, since they have useful properties with the mentioned adjunction. Then, I will introduce the notion of nullary and essentially nullary monads and their results, which will be used the prove the main result of the ideally exact categories.

In the third chapter I will present some results about the Eilenberg-Moore category of a monad. This part is quite technical and articulate. It will be useful to characterize the notion of ideal exact category. After these results, I will present the definition of ideal exact category and some of its properties. Here I will also prove the main result of the ideally exact category, which has been central in motivating interest in this topic.